\paragrafos{Απογραφή, πληθυσμός, κατηγορικά και ποσοτικά δεδομένα}{Γ.1.1, «Απογραφή – Πληθυσμός – Κατηγορικά δεδομένα – Ποσοτικά δεδομένα»}

\begin{anagke}
Η κυρία Αλεξάνδρα ανακοινώνει ότι το σχολείο θα οργανώσει αθλητική ημέρα, αλλά το κονδύλι φτάνει για ένα μόνο άθλημα. «Ποιο να διαλέξουμε;» ρωτά. Ο Νίκος φωνάζει αμέσως «μπάσκετ!». Η Ζωή απαντά: «Εσύ το θέλεις, αλλά το σχολείο;». Ο Στέλιος προτείνει να ρωτήσουν μόνο τους φίλους του, γιατί «είναι πιο γρήγορο». Η Ρήνη αντιδρά: «Οι φίλοι σου παίζουν όλοι ποδόσφαιρο, δεν θα βγει νόημα».

Η Φανή πιάνει το τετράδιό της και προτείνει: «Να ρωτήσουμε τους συμμαθητές μας, να σημειώσουμε τις απαντήσεις και να δούμε τι βγαίνει». Αμέσως προκύπτουν ερωτήματα: ποιους θα ρωτήσουμε, τι ακριβώς θα ρωτήσουμε, πώς θα γράψουμε τις απαντήσεις ώστε να διαβάζονται με μια ματιά.

Σε αυτή την παράγραφο μαθαίνουμε να ονομάζουμε αυτό που ερευνούμε (πληθυσμός, μεταβλητή), να διαλέγουμε σωστά ποιους θα ρωτήσουμε (δείγμα), να ξεχωρίζουμε τα είδη των δεδομένων και να τα οργανώνουμε σε πίνακα συχνοτήτων και ραβδόγραμμα.
\end{anagke}

\begin{lexi}{στατιστική, απογραφή, πληθυσμός, δείγμα, μεταβλητή}{ονόματα της στατιστικής}
Η \lexiword{στατιστική} πήρε το όνομά της από το «κράτος»: οι πρώτοι που μάζευαν δεδομένα ήταν οι κυβερνήσεις, για να ξέρουν πόσοι είναι οι πολίτες τους. Η \lexiword{απογραφή} είναι «το γράψιμο από την αρχή ως το τέλος» όλων. Ο \lexiword{πληθυσμός} είναι ο «πολύς κόσμος» που μας ενδιαφέρει. Το \lexiword{δείγμα} είναι το «παράδειγμα»: ένα μικρό κομμάτι που δείχνει το όλο, όπως η γεύση ενός κουταλιού από το φαγητό. Η \lexiword{μεταβλητή} είναι αυτό που αλλάζει από άτομο σε άτομο.
\end{lexi}

\begin{metafrasi}
Όλοι όσοι μας ενδιαφέρουν. & πληθυσμός\\
Ρωτάω όλους, έναν έναν. & απογραφή\\
Ρωτάω μόνο μερικούς, επιλεγμένους με τύχη. & δείγμα\\
Αυτό που ρωτάω (ύψος, χρώμα, άθλημα). & μεταβλητή\\
Οι απαντήσεις που κατέγραψα. & στατιστικά δεδομένα\\
Πόσες φορές εμφανίστηκε μια τιμή. & συχνότητα\\
Ποιο μέρος του συνόλου είναι μια τιμή. & σχετική συχνότητα\\
\end{metafrasi}

\tmima{Η θεωρία}

\begin{orismos}{πληθυσμός και μεταβλητή}
\textbf{Πληθυσμός} είναι μια ομάδα ανθρώπων ή ένα σύνολο αντικειμένων, των οποίων τα στοιχεία μελετάμε ως προς κάποιο χαρακτηριστικό τους. Αυτό το χαρακτηριστικό ονομάζεται \textbf{μεταβλητή}.
\end{orismos}

Παραδείγματα πληθυσμών: οι μαθητές των γυμνασίων της χώρας, τα αυτοκίνητα που κινούνται σε μια εθνική οδό, τα πρόβατα του παππού Γιάννη. Παραδείγματα μεταβλητών: το ύψος των μαθητών, το χρώμα ενός αυτοκινήτου, το βάρος ενός προβάτου.

\begin{orismos}{απογραφή, δείγμα, δειγματοληψία}
Η διαδικασία με την οποία εξετάζουμε και καταγράφουμε \emph{όλα} τα στοιχεία ενός πληθυσμού λέγεται \textbf{απογραφή}. Όταν η απογραφή κοστίζει πολύ σε χρόνο ή σε χρήματα, εξετάζουμε μόνο ένα μέρος του πληθυσμού, που λέγεται \textbf{δείγμα}. Η διαδικασία επιλογής του δείγματος λέγεται \textbf{δειγματοληψία}.
\end{orismos}

Το δείγμα πρέπει να επιλέγεται έτσι ώστε τα συμπεράσματα που θα βγάλουμε να είναι \textbf{αξιόπιστα}. Τότε λέμε ότι το δείγμα είναι \textbf{αντιπροσωπευτικό} του πληθυσμού: μοιάζει με τον πληθυσμό, σαν μικρογραφία του. Αν ρωτήσουμε τον πληθυσμό ή το δείγμα, οι απαντήσεις που συγκεντρώνουμε λέγονται \textbf{στατιστικά δεδομένα} ή \textbf{παρατηρήσεις}.

\begin{kanonas}{τα είδη των δεδομένων}
\textbf{Κατηγορικά} είναι τα δεδομένα που δεν είναι αριθμοί (χρώμα ματιών, άθλημα, μάρκα κινητού).\\
\textbf{Ποσοτικά} είναι τα δεδομένα που είναι αριθμοί (ύψος, ηλικία, αριθμός αδελφών). Τα ποσοτικά χωρίζονται σε δύο είδη:\\
\textbf{διακριτά}, όταν παίρνουν μεμονωμένες τιμές (π.χ. αριθμός αδελφών: $0$, $1$, $2$, $3$ και όχι $1{,}5$)\\
\textbf{συνεχή}, όταν μπορούν να πάρουν οποιαδήποτε τιμή μέσα σε ένα διάστημα (π.χ. ο χρόνος σε λεπτά που χρειάζεται ένας μαθητής για να λύσει ένα διαγώνισμα, από $0$ ως $120$).
\end{kanonas}

\begin{center}
\begin{tikzpicture}[font=\small,
  nd/.style={draw,rounded corners=3pt,align=center,inner sep=4pt,fill=gblueL},
  nl/.style={draw,rounded corners=3pt,align=center,inner sep=4pt,fill=gwarm}]
\node[nd] (a) at (0,0) {Στατιστικά\\ δεδομένα};
\node[nd] (b) at (4,1.1) {Ποσοτικά\\ (είναι αριθμοί)};
\node[nd] (c) at (4,-1.3) {Κατηγορικά\\ (δεν είναι αριθμοί)};
\node[nl] (d) at (8.3,1.9) {Διακριτά\\ (μεμονωμένες τιμές)};
\node[nl] (e) at (8.3,0.3) {Συνεχή\\ (κάθε τιμή σε διάστημα)};
\draw (a.east)--++(0.6,0)|-(b.west);
\draw (a.east)--++(0.6,0)|-(c.west);
\draw (b.east)--++(0.6,0)|-(d.west);
\draw (b.east)--++(0.6,0)|-(e.west);
\end{tikzpicture}
\end{center}

\begin{prosoxi}
\begin{itemize}[leftmargin=1.2em,itemsep=0.3ex,topsep=2pt]
\item Δεν αρκεί να είναι κάτι γραμμένο με αριθμό για να είναι ποσοτικό. Ο αριθμός της φανέλας ενός παίκτη δεν μετρά τίποτα, απλώς τον ξεχωρίζει: είναι κατηγορικό δεδομένο. Ποσοτικό είναι ό,τι μετράει μια ποσότητα, δηλαδή απαντά στο «πόσο» ή στο «πόσα».
\item Αν το δείγμα έχει επιλεγεί από ένα μόνο τμήμα του πληθυσμού (π.χ. μόνο οι φίλοι μου, μόνο η ομάδα μπάσκετ), τα συμπεράσματα δεν είναι αξιόπιστα για όλο τον πληθυσμό.
\item Το πλήθος (αριθμός παιδιών, αριθμός αδελφών) είναι διακριτό. Το ύψος, το βάρος και ο χρόνος είναι συνεχή.
\end{itemize}
\end{prosoxi}

\begin{orismos}{συχνότητα και σχετική συχνότητα}
\textbf{Συχνότητα} μιας τιμής είναι το πλήθος των φορών που εμφανίζεται στα δεδομένα. \textbf{Σχετική συχνότητα} μιας τιμής είναι το πηλίκο της συχνότητάς της διά του συνόλου των παρατηρήσεων. Τη γράφουμε συνήθως ως ποσοστό επί τοις εκατό ($\%$).
\end{orismos}

Οι συχνότητες και οι σχετικές συχνότητες συγκεντρώνονται σε έναν \textbf{πίνακα συχνοτήτων}. Το άθροισμα όλων των συχνοτήτων είναι το σύνολο των παρατηρήσεων και το άθροισμα όλων των σχετικών συχνοτήτων είναι πάντα $100\%$ (ή $1$).

Το \textbf{ραβδόγραμμα} παριστάνει τις συχνότητες με ράβδους ίδιου πλάτους. Το ύψος κάθε ράβδου ισούται με τη συχνότητα της τιμής και οι ράβδοι δεν ακουμπούν μεταξύ τους. Είναι κατάλληλο για κατηγορικά και για διακριτά δεδομένα.

\tmima{Λυμένα παραδείγματα}

\paradeigma{Ποιος είναι ο πληθυσμός;}
\tixero Η Ρήνη θέλει να μάθει πόσα λεπτά την ημέρα παίζουν παιχνίδια στο κινητό οι $280$ μαθητές του σχολείου. Ρωτά τυχαία $30$ μαθητές από όλες τις τάξεις.\par
\tithelo Πληθυσμός, δείγμα, μεταβλητή και είδος δεδομένων.\par
\possindeo Εξετάζουμε όλους τους μαθητές του σχολείου: αυτοί είναι ο πληθυσμός. Ρωτήθηκαν μόνο $30$: αυτοί είναι το δείγμα. Αυτό που ρωτάμε είναι ο χρόνος παιχνιδιού.\par
\apantisi Πληθυσμός: οι $280$ μαθητές. Δείγμα: οι $30$ μαθητές. Μεταβλητή: τα λεπτά παιχνιδιού την ημέρα. Τα δεδομένα είναι ποσοτικά και συνεχή, γιατί ο χρόνος μπορεί να πάρει οποιαδήποτε τιμή (π.χ. $42{,}5$ λεπτά). Το δείγμα είναι αντιπροσωπευτικό, αφού τα παιδιά επελέγησαν τυχαία από όλες τις τάξεις.\par
\elegxos Το δείγμα ($30$) είναι μικρότερο από τον πληθυσμό ($280$), όπως πρέπει.

\paradeigma{Αντιπροσωπευτικό ή όχι}
\tixero Ο Ηλίας θέλει να μάθει το ύψος των μαθητών του σχολείου και μετρά τα $22$ παιδιά της ομάδας μπάσκετ της Γ΄ τάξης.\par
\tithelo Αν το δείγμα είναι αντιπροσωπευτικό.\par
\possindeo Οι παίκτες του μπάσκετ είναι συνήθως ψηλότεροι. Επιπλέον είναι μόνο της Γ΄ τάξης, που έχει μεγαλύτερα παιδιά από την Α΄ και τη Β΄.\par
\apantisi Δεν είναι αντιπροσωπευτικό. Τα συμπεράσματα δεν θα είναι αξιόπιστα για όλο το σχολείο.\par
\elegxos Αν είχε επιλέξει $22$ παιδιά με κλήρωση από όλες τις τάξεις, το δείγμα θα ήταν καλύτερο.

\paradeigma{Είδος δεδομένων}
\tixero Οι μεταβλητές: χρώμα ματιών, αριθμός αδελφών, βάρος σχολικής τσάντας.\par
\tithelo Το είδος των δεδομένων της καθεμιάς.\par
\possindeo Το χρώμα δεν είναι αριθμός. Ο αριθμός αδελφών είναι αριθμός με μεμονωμένες τιμές. Το βάρος είναι αριθμός και μπορεί να πάρει οποιαδήποτε τιμή (π.χ. $3{,}27$ kg).\par
\apantisi Χρώμα ματιών: κατηγορικά. Αριθμός αδελφών: ποσοτικά διακριτά. Βάρος τσάντας: ποσοτικά συνεχή.\par
\elegxos Το «πόσα» δίνει διακριτά, το «πόσο» δίνει συνεχή.

\paradeigma{Πίνακας συχνοτήτων και ραβδόγραμμα}
\tixero Η Φανή ρώτησε και τα $25$ παιδιά της Α$3$ ποιο είναι το αγαπημένο τους άθλημα. Ποδόσφαιρο είπαν $8$, μπάσκετ $6$, βόλεϊ $4$, κολύμβηση $3$, τένις $2$ και άλλο $2$.\par
\tithelo Τον πίνακα συχνοτήτων με τις σχετικές συχνότητες και το ραβδόγραμμα.\par
\possindeo Κάθε σχετική συχνότητα είναι η συχνότητα διά $25$: $\frac{8}{25}=0{,}32=32\%$, $\frac{6}{25}=0{,}24=24\%$, και ούτω καθεξής.\par
\apantisi

\begin{center}
\renewcommand{\arraystretch}{1.25}
\begin{tabular}{lcc}
\toprule
\textbf{Άθλημα} & \textbf{Συχνότητα} & \textbf{Σχετική συχνότητα}\\
\midrule
Ποδόσφαιρο & $8$ & $32\%$\\
Μπάσκετ & $6$ & $24\%$\\
Βόλεϊ & $4$ & $16\%$\\
Κολύμβηση & $3$ & $12\%$\\
Τένις & $2$ & $8\%$\\
Άλλο & $2$ & $8\%$\\
\midrule
\textbf{Σύνολο} & $25$ & $100\%$\\
\bottomrule
\end{tabular}
\end{center}

\begin{center}
\gaRabdos{10}{2}{Ποδόσφαιρο/8,Μπάσκετ/6,Βόλεϊ/4,Κολύμβηση/3,Τένις/2,Άλλο/2}{παιδιά}
\end{center}

\elegxos $8+6+4+3+2+2=25$ και $32+24+16+12+8+8=100$.

\tmima{Ασκήσεις}

\zoni{Το κατάλαβα}{ονόματα και είδη}

\ex Κατηγορικά ή ποσοτικά δεδομένα; Στα ποσοτικά πες αν είναι διακριτά ή συνεχή. α) Το χρώμα της φανέλας μιας ομάδας. β) Ο αριθμός των αδελφών ενός παιδιού. γ) Το αγαπημένο γλυκό. δ) Οι ώρες ύπνου μιας νύχτας. ε) Το βάρος μιας σχολικής τσάντας. στ) Ο αριθμός των επαφών στο κινητό. ζ) Η μάρκα ενός κινητού.

\ex Η Ρήνη θέλει να μάθει πόσα βιβλία διάβασαν το καλοκαίρι οι $120$ μαθητές της Α΄ τάξης του σχολείου της. Ρωτά $20$ παιδιά που διάλεξε με κλήρωση. Ποιος είναι ο πληθυσμός, ποιο το δείγμα, ποια η μεταβλητή και τι είδους είναι τα δεδομένα;

\ex Απογραφή ή δείγμα; α) Η κυρία Αλεξάνδρα καταγράφει τις απουσίες και των $25$ παιδιών της Α$3$. β) Ο δήμος ρωτά $500$ κατοίκους από τους $12\,000$ για το νέο πάρκο. γ) Ο παππούς Γιάννης δοκιμάζει ένα σταφύλι για να κρίνει αν ο αμπελώνας είναι γλυκός. δ) Ο Ηλίας μετρά το ύψος όλων των παιδιών της τάξης του.

\ex Η Αγγελική κατέγραψε τι έφεραν για κολατσιό τα $20$ παιδιά μιας τάξης. Τοστ έφεραν $7$, φρούτο $5$, κουλούρι $4$, γιαούρτι $2$ και τίποτα $2$ παιδιά. Φτιάξε πίνακα συχνοτήτων με σύνολο και σχετικές συχνότητες.

\ex Στο ραβδόγραμμα φαίνεται πώς πηγαίνουν στο σχολείο τα $25$ παιδιά της Α$3$.
\begin{center}
\gaRabdos{10}{2}{Με τα πόδια/9,Ποδήλατο/3,Λεωφορείο/5,Αυτοκίνητο/8}{παιδιά}
\end{center}
α) Ποιος τρόπος είναι ο πιο συνηθισμένος και ποιος ο πιο σπάνιος; β) Πόσα παιδιά δεν έρχονται με αυτοκίνητο; γ) Ποια είναι η σχετική συχνότητα του ποδηλάτου;

\ex Η γιαγιά Μυρτώ έχει στο χωριό $12$ κότες, $8$ πρόβατα, $6$ γίδες, $4$ γάτες και $2$ σκύλους. Φτιάξε πίνακα συχνοτήτων και ραβδόγραμμα. Πόσα ζώα έχει συνολικά;

\ex Σωστό ή λάθος; α) Ο πληθυσμός έχει πάντα ανθρώπους. β) Το δείγμα έχει συνήθως λιγότερα στοιχεία από τον πληθυσμό. γ) Ο αριθμός των αδελφών είναι διακριτό δεδομένο. δ) Το χρώμα ματιών είναι ποσοτικό δεδομένο. ε) Στο ραβδόγραμμα το ύψος της ράβδου δείχνει τη συχνότητα.

\zoni{Το χρησιμοποιώ}{ερωτηματολόγια και πίνακες}

\ex \textbf{Το δείγμα.} Ο Στέλιος θέλει να μάθει ποιο τραγούδι προτιμούν οι μαθητές του σχολείου. Προτείνει τρεις τρόπους: α) να ρωτήσει τους $10$ φίλους του, β) να ρωτήσει $30$ παιδιά που θα διαλέξει με κλήρωση από όλες τις τάξεις, γ) να ρωτήσει όσους βρει στο κυλικείο το πρώτο διάλειμμα. Ποιος τρόπος δίνει το πιο αντιπροσωπευτικό δείγμα και γιατί;

\ex \textbf{Τα αυτοκίνητα.} Η Φανή κάθεται έξω από το σχολείο και καταγράφει το χρώμα $20$ αυτοκινήτων που περνούν: Λευκό (Α), Γκρι (Γ), Μαύρο (Μ), Κόκκινο (Κ), Άλλο (Ο).\par
\noindent Γ Κ Ο Μ Γ Μ Α Α Γ Α Μ Α Α Γ Κ Α Μ Α Α Γ\par
Φτιάξε πίνακα συχνοτήτων με σχετικές συχνότητες και ραβδόγραμμα.

\ex \textbf{Τα βιβλία.} Στο ραβδόγραμμα φαίνεται πόσα βιβλία διάβασαν το καλοκαίρι τα $30$ παιδιά μιας τάξης.
\begin{center}
\gaRabdos{10}{2}{0 βιβλία/4,1 βιβλίο/9,2 βιβλία/8,3 βιβλία/6,4 βιβλία/3}{παιδιά}
\end{center}
α) Ποια τιμή έχει τη μεγαλύτερη συχνότητα; β) Πόσα παιδιά διάβασαν τουλάχιστον $2$ βιβλία; γ) Ποια είναι η σχετική συχνότητα των παιδιών που διάβασαν ένα βιβλίο; δ) Πόσα βιβλία διάβασαν συνολικά όλα τα παιδιά;

\ex \textbf{Τα κινητά.} Ο πίνακας δείχνει τέσσερα κινητά ενός καταστήματος.
\begin{center}
\renewcommand{\arraystretch}{1.2}
\begin{tabular}{lcccc}
\toprule
\textbf{Μοντέλο} & \textbf{Χρώμα} & \textbf{Μνήμη (GB)} & \textbf{Βάρος (g)} & \textbf{Αδιάβροχο}\\
\midrule
Κ1 & μπλε & $64$ & $172{,}5$ & ναι\\
Κ2 & μαύρο & $128$ & $180$ & ναι\\
Κ3 & λευκό & $64$ & $165{,}2$ & όχι\\
Κ4 & μπλε & $256$ & $190{,}8$ & ναι\\
\bottomrule
\end{tabular}
\end{center}
Για καθεμιά από τις στήλες από τη δεύτερη ως την πέμπτη, πες αν τα δεδομένα είναι κατηγορικά ή ποσοτικά και, αν είναι ποσοτικά, διακριτά ή συνεχή.

\ex \textbf{Ο συμπληρωμένος πίνακας.} Στο ραβδόγραμμα της Ρήνης φαίνονται οι δραστηριότητες που παρακολουθούν μετά το σχολείο $40$ παιδιά. Οι ράβδοι της μουσικής, του χορού, της ζωγραφικής και της ρομποτικής έχουν ύψη $10$, $8$, $6$ και $12$. Η πέμπτη ράβδος είναι του σκακιού και δεν φαίνεται στο χαρτί. Πόσο είναι το ύψος της και ποια η σχετική συχνότητα κάθε δραστηριότητας;

\zoni{Το σκέφτομαι}{γιατί και πώς}

\ex Ο Βασίλης λέει: «Από την έρευνα της Φανής, $8$ στα $25$ παιδιά της Α$3$ προτιμούν το ποδόσφαιρο, άρα το $32\%$ όλων των μαθητών του σχολείου προτιμά το ποδόσφαιρο». Συμφωνείς; Τι θα έπρεπε να ισχύει για το δείγμα;

\ex Η Ζωή θέλει να μάθει πόσες ώρες κοιμούνται τα παιδιά της ηλικίας της. Όρισε: πληθυσμό, δείγμα, μεταβλητή, είδος δεδομένων, και πρότεινε πώς θα επιλεγεί το δείγμα.

\ex Η Αγγελική βρήκε τις σχετικές συχνότητες $40\%$, $35\%$ και $30\%$ για τρεις κατηγορίες. Ο Ηλίας λέει ότι κάτι δεν πάει καλά. Ποιος έχει δίκιο και γιατί;

\ex Σκέψου μια ερώτηση που μπορείς να κάνεις στην Α$3$ και θα δώσει: α) κατηγορικά δεδομένα, β) διακριτά ποσοτικά δεδομένα, γ) συνεχή ποσοτικά δεδομένα.

\zoni{Για το σπίτι}{ερευνώ στην οικογένεια}

\ex Ρώτησε $10$ άτομα της οικογένειας ή της γειτονιάς ποιο είναι το αγαπημένο τους φαγητό. Φτιάξε πίνακα συχνοτήτων και ραβδόγραμμα.

\ex Βρες σε εφημερίδα, περιοδικό ή ιστοσελίδα ένα διάγραμμα. Γράψε ποιος είναι ο πληθυσμός, ποια η μεταβλητή και τι είδους είναι τα δεδομένα.
